\chapter{Ordered pairs, products, and relations}\label{ch:relations}
\begin{readergoal}Understand how ordered observations lead to relations and why correlation, equality, and ordering are mathematically different relations.\end{readergoal}
\section{The order of information matters}
The collection $\{2,7\}$ does not specify which element came first. An \emph{ordered pair} $(2,7)$ does. Formally, an ordered-pair construction must satisfy
\[(a,b)=(c,d)\quad\Longleftrightarrow\quad a=c\ \text{and}\ b=d.\]
One set-theoretic realization is $(a,b)=\{\{a\},\{a,b\}\}$, although the particular encoding is much less important than the stated property. A spike observation might be the pair $(t,u)$ of timestamp and unit identifier; reversing it gives a different type of object.

The \emph{Cartesian product} $A\times B$ contains all ordered pairs $(a,b)$ with $a\in A$ and $b\in B$. If $A$ has $m$ elements and $B$ has $n$, then $A\times B$ has $mn$ elements. More generally, $A_1\times\cdots\times A_k$ defines ordered tuples. A multichannel recording may be indexed by channel $c$, time sample $n$, and trial $r$: its index space is $C\times T\times R$.

\section{What is a relation?}
A \emph{binary relation} from $A$ to $B$ is a subset $R\subseteq A\times B$. Writing $aRb$ is shorthand for $(a,b)\in R$. One relation on the integers is $a\le b$; another is ``$a$ divides $b$.'' In neuroscience, an empirically defined relation may connect two neurons if their estimated spike-count correlation is above a threshold. This makes a graph edge, not necessarily a synapse.

A relation $R\subseteq A\times A$ is \emph{reflexive} if $aRa$ for every $a$, \emph{symmetric} if $aRb$ implies $bRa$, and \emph{transitive} if $aRb$ and $bRc$ imply $aRc$. A relation satisfying all three is an \emph{equivalence relation}; it partitions $A$ into disjoint equivalence classes. A \emph{partial order} is reflexive, antisymmetric, and transitive. Equality is an equivalence relation; subset inclusion is a partial order; temporal precedence is an order relation that requires a convention for simultaneous events.

\section{Composing and inverting relations}
For $R\subseteq A\times B$ and $S\subseteq B\times C$, define
\[S\circ R=\{(a,c):\exists b\in B,\ (a,b)\in R\ \land\ (b,c)\in S\}.\]
The inverse relation is $R^{-1}=\{(b,a):(a,b)\in R\}$. Relation inversion is not yet the same as inversion of a function. For instance, a relation recording which neurons are active during which trials can be reversed to ask which trials contain activity from a neuron.

\begin{workedexample}[Why graph clustering needs care]
Let $uRv$ mean ``estimated Pearson correlation between $u$ and $v$ exceeds 0.5.'' The rule is symmetric whenever correlations are defined, and self-pairs have correlation one for nonconstant variables, but it need not be transitive: neuron 1 may correlate strongly with 2, 2 with 3, and 1 weakly with 3. Thus its connected components are not the same as equivalence classes generated by the direct threshold relation.
\end{workedexample}
% --- expanded edition (October 2026) ---
\section{Equivalence relations and partitions}
Equivalence relations and partitions are two descriptions of the same structure. Given an equivalence relation $\sim$ on $A$, the \emph{equivalence class} of $a$ is $[a]=\{x\in A:x\sim a\}$, and the \emph{quotient set} is $A/{\sim}=\{[a]:a\in A\}$.

\begin{theorem}[Fundamental theorem of equivalence relations]
Let $\sim$ be an equivalence relation on $A$. Then
\begin{enumerate}[leftmargin=2em]
\item every $a\in A$ belongs to its own class $[a]$;
\item for $a,b\in A$, either $[a]=[b]$ or $[a]\cap[b]=\varnothing$;
\item consequently the classes form a partition of $A$.
\end{enumerate}
Conversely, every partition $\mathcal P$ of $A$ defines an equivalence relation by $a\sim b$ iff $a$ and $b$ lie in the same block of $\mathcal P$, and its classes are exactly the blocks.
\end{theorem}
\begin{proof}
(1) follows from reflexivity. For (2), suppose $c\in[a]\cap[b]$, so $c\sim a$ and $c\sim b$. If $x\in[a]$ then $x\sim a$; by symmetry $a\sim c$; by transitivity $x\sim c$ and then $x\sim b$, so $x\in[b]$. Hence $[a]\subseteq[b]$, and the reverse inclusion follows by exchanging the roles of $a$ and $b$. (3) The classes are nonempty by (1), cover $A$ by (1), and are pairwise disjoint or equal by (2). For the converse, reflexivity and symmetry are immediate, and transitivity holds because the blocks are disjoint: if $a,b$ share a block and $b,c$ share a block, both blocks contain $b$ and therefore coincide.
\end{proof}

This theorem underlies all the number-system constructions of Part~II: integers and rationals are equivalence classes of pairs. It also explains why ``spike sorting'' is mathematically a partition problem: an assignment of each spike to exactly one unit is the same as an equivalence relation ``was fired by the same unit as.''

\section{Closures of relations}
A relation that fails to be transitive can be enlarged to a transitive one. The \emph{transitive closure} $R^+$ of $R\subseteq A\times A$ is the smallest transitive relation containing $R$. Concretely, $aR^+b$ iff there is a finite chain $a=x_0Rx_1R\cdots Rx_k=b$ with $k\ge1$. The \emph{reflexive--symmetric--transitive closure} of $R$ is the smallest equivalence relation containing $R$; its classes are the \emph{connected components} of the undirected graph whose edges are the pairs in $R$.

\begin{proposition}
The intersection of any nonempty family of transitive relations on $A$ is transitive. Hence the transitive closure exists and equals the intersection of all transitive relations containing $R$.
\end{proposition}
\begin{proof}
Let $(T_i)_{i\in I}$ be transitive and $T=\bigcap_iT_i$. If $aTb$ and $bTc$, then for each $i$ we have $aT_ib$ and $bT_ic$, hence $aT_ic$. Thus $aTc$. The family of transitive relations containing $R$ is nonempty, since $A\times A$ belongs to it, so its intersection is transitive, contains $R$, and is contained in every transitive relation containing $R$.
\end{proof}

The worked example on correlation thresholds above can now be stated precisely: clustering by connected components replaces the threshold relation by its equivalence closure, which may merge neurons that are not directly correlated at all. Long ``chains'' of moderately correlated units can join two unrelated groups.

\section{Orders}
A \emph{partial order} $\preceq$ on $A$ is reflexive, antisymmetric ($a\preceq b$ and $b\preceq a$ imply $a=b$), and transitive. It is a \emph{total order} if, in addition, any two elements are comparable. The usual $\le$ on numbers is total; inclusion $\subseteq$ on the power set of a set with at least two elements is only partial, since $\{u\}$ and $\{v\}$ are incomparable. A \emph{strict} order $<$ is irreflexive and transitive; from a partial order one obtains a strict one by $a<b\iff a\preceq b\land a\ne b$.

Orders appear in data whenever there is ranking: spike times are totally ordered, but the relation ``neuron $u$ tends to fire before neuron $v$'' estimated from cross-correlograms need not even be transitive. Sequence-detection methods that assume a consistent total order of activation must test, not presuppose, that the estimated precedence relation has that structure.

\begin{workedexample}[Modular phase]
On the real line define $x\sim y$ iff $x-y=2\pi k$ for some integer $k$. This is an equivalence relation: $x-x=0$; if $x-y=2\pi k$ then $y-x=2\pi(-k)$; and $x-y=2\pi k$, $y-z=2\pi m$ give $x-z=2\pi(k+m)$. The quotient $\mathbb R/{\sim}$ is the circle of phases. Each class has exactly one representative in $(-\pi,\pi]$, which is the ``wrapped'' phase returned by \texttt{angle}. Arithmetic on representatives is well defined for addition (adding phases) but not for ordering: ``phase $\alpha$ is less than phase $\beta$'' depends on the representative chosen.
\end{workedexample}
\begin{exercises}
\item List $A\times B$ for $A=\{a,b\}$, $B=\{0,1\}$, and explain why $A\times B\ne B\times A$ generally.
\item Construct a reflexive symmetric relation that is not transitive.
\item Prove the inverse of an equivalence relation is the same relation.
\item Interpret a directed graph edge as a relation; explain why its direction does not prove causal direction.
\item Show that ``$a\equiv b \pmod 3$'' is an equivalence relation on $\mathbb Z$ and list its classes.
\item Compute the transitive closure of $R=\{(1,2),(2,3),(3,4)\}$ on $\{1,2,3,4\}$.
\item Prove that the inverse of a partial order is a partial order.
\item Give a precedence relation among three neurons that is not transitive, and explain why no total order is consistent with it.
\item Let $A$ be a finite set with a partition into $k$ blocks. How many equivalence relations on $A$ have exactly these classes?
\end{exercises}
